\subsubsection{Class: IoTSensor}
\textbf{Mô tả (Overview):} Lớp đại diện cho các thiết bị phần cứng cảm biến (ví dụ: cảm biến siêu âm, từ trường) được lắp đặt tại mỗi vị trí đỗ xe để phát hiện sự hiện diện của phương tiện.

\vspace{0.3cm}
\noindent\textbf{Thuộc tính (Attributes):}
\begin{itemize}
    \item \texttt{- sensorId: String}: Định danh duy nhất của cảm biến (UUID hoặc MAC Address).
    \item \texttt{- slotId: String}: Khóa ngoại (Foreign Key) liên kết với \texttt{ParkingSlot} mà cảm biến này đang giám sát.
    \item \texttt{- isActive: Boolean}: Trạng thái hoạt động của cảm biến (\texttt{True} nếu đang online và hoạt động tốt, \texttt{False} nếu mất kết nối hoặc hỏng).
\end{itemize}

\noindent\textbf{Phương thức (Methods):}
\begin{itemize}
    \item \texttt{+ transmitOccupancyData(): void}: Kích hoạt khi cảm biến phát hiện sự thay đổi trạng thái vật lý (từ trống sang có xe hoặc ngược lại). Đóng gói dữ liệu và gửi payload (qua MQTT/HTTP) về IoT Gateway.
    \item \texttt{+ checkHealth(): void}: Gửi tín hiệu nhịp tim (heartbeat) định kỳ về server chứa các thông số như \texttt{nguồn điện} và tình trạng mạng để hệ thống biết cảm biến vẫn đang hoạt động (Online).
\end{itemize}


\subsubsection{Class: ElectronicSignage}
\textbf{Mô tả (Overview):} Lớp đại diện cho các bảng hiển thị điện tử (LED/LCD) được lắp đặt tại các điểm nút giao thông trong trường hoặc trước cổng bãi xe nhằm điều hướng và thông báo số chỗ trống.

\vspace{0.3cm}
\noindent\textbf{Thuộc tính (Attributes):}
\begin{itemize}
    \item \texttt{- signageId: String}: Định danh duy nhất của bảng hiển thị.
    \item \texttt{- location: String}: Vị trí vật lý lắp đặt bảng (Ví dụ: "Cổng 1", "Ngã tư Khu A", "Tầng hầm B1").
    \item \texttt{- displayMessage: String}: Nội dung văn bản hoặc thông điệp hiện tại đang được hiển thị trên bảng.
\end{itemize}

\noindent\textbf{Phương thức (Methods):}
\begin{itemize}
    \item \texttt{+ updateMessage(text: String): void}: Phương thức nhận lệnh từ hệ thống trung tâm (\texttt{ParkingManagementService}) để cập nhật lại biến \texttt{displayMessage} và render dòng chữ mới ra màn hình LED (Ví dụ: "Khu A: Còn 5 chỗ", "Khu B: HẾT CHỖ").
\end{itemize}